\hiddenappendixchapter
    {Evolution of the Inherited Inspection Modules}
    {app:inherited_interface_evolution}
% Compatibility with the reference used by the previous appendix version.
\label{app:application_interface_overview}

\providecommand{\appAimage}[1]{%
    \IfFileExists{#1}{%
        \includegraphics[
            width=\linewidth,
            height=0.30\textheight,
            keepaspectratio
        ]{#1}%
    }{%
        \fbox{%
            \parbox[c][0.19\textheight][c]{0.94\linewidth}{%
                \centering
                \textbf{Capture pending}\\[0.6em]
                \texttt{\detokenize{#1}}
            }%
        }%
    }%
}

The inspection modules retained from the initial application were adapted to
use the common case collection and model hierarchy described in
Section~\ref{sec:final_application_interface}. Their main purpose remains the
inspection of groups, nodes and thermal connections at a selected model state,
but the final implementation applies the same physical selection to several
compatible cases and evaluates transient data at a specified output instant.

The comparisons below focus on the interface changes rather than on new
thermal results. The nominal case and the same hierarchy spreadsheet should be
loaded in both versions. The hot and cold cases are added only in the final
interface when required to expose the multi-case controls. Browser width,
zoom, selected groups and displayed units should remain unchanged between each
pair of captures.

\section{Flow Explorer}
\label{appsec:flow_explorer_evolution}

The initial Flow Explorer provides hierarchical source and destination
selections and generates Sankey, network and bar representations for the
active dataset. Its group lists are assembled from the available hierarchy,
but the initial interface does not consistently preserve the parent--child
order and permits redundant selections in which a parent and one of its
descendants are included simultaneously.

The final module orders the available groups according to the Excel hierarchy
and removes descendant selections already contained in a selected parent. A
visible message identifies the corrected selection. The same controls are
then applied to one or more cases at a selected output time. Superposed and
\textit{Case difference} representations are available, and the generated
figure can be exported or stored together with the configuration required for
later restoration.

% Capture with the same root, source, destination, units and graphical type.
% Recommended selection: root AVIONICS, source OBC_BOX, destination RADIATOR,
% total heat flow, t = 0 s. In the final version also load the hot and cold
% cases and keep the multi-case controls visible.
\begin{figure}[!htbp]
    \centering

    \begin{subfigure}[t]{0.9\textwidth}
        \centering
        \appAimage{Figures/Anexo_01/original_flow_explorer.png}
        \caption{Initial single-case interface.}
        \label{fig:appendix_original_flow_explorer}
    \end{subfigure}

    \vspace{0.8em}

    \begin{subfigure}[t]{0.96\textwidth}
        \centering
        \appAimage{Figures/Anexo_01/final_flow_explorer.png}
        \caption{Final transient and multi-case interface.}
        \label{fig:appendix_final_flow_explorer}
    \end{subfigure}

    \caption{Evolution of the Flow Explorer interface.}
    \label{fig:appendix_flow_explorer_evolution}
\end{figure}

\FloatBarrier

\section{Case Information}
\label{appsec:case_information_evolution}

The initial \textit{Model info} tab reports the node, group and coupling data
associated with the active result file. The final \textit{Case info} tab
retains these summaries for every selected case and adds level-based coupling
inspection at a chosen output time. The visible case name is kept separate
from the internal analysis identifier, allowing files with the same internal
name to remain distinguishable throughout the interface.

% Original: show the main model metrics and one generated group-information
% block. Final: show the corresponding Case info metrics for the nominal case,
% plus the level coupling selector. Crop both images to the active tab.
\begin{figure}[!htbp]
    \centering

    \begin{subfigure}[t]{0.96\textwidth}
        \centering
        \appAimage{Figures/Anexo_01/original_model_info.png}
        \caption{Initial model-information view.}
        \label{fig:appendix_original_model_info}
    \end{subfigure}

    \vspace{0.8em}

    \begin{subfigure}[t]{0.96\textwidth}
        \centering
        \appAimage{Figures/Anexo_01/final_case_info.png}
        \caption{Final case-information view.}
        \label{fig:appendix_final_case_info}
    \end{subfigure}

    \caption{Evolution from active-model information to case-specific
    inspection.}
    \label{fig:appendix_case_information_evolution}
\end{figure}

\FloatBarrier

\section{Group Attributes}
\label{appsec:group_attributes_evolution}

The initial group table is generated for one active case and reads the first
available temporal row. This behaviour is sufficient for a steady-state
solution but does not allow the inspected instant to be selected in a
transient dataset. The final tab uses the stored output-time vector, evaluates
the selected attributes at the requested instant and creates independent
tables for the active cases. The same area can also calculate reference-minus-
comparison tables when \textit{Case difference} is selected.

% Use BATTERY, OBC_BOX and RADIATOR with Tmean, QI and QR. Generate the table
% at t = 0 s in both versions. In the final capture select the three cases and
% keep the representation selector visible above the table.
\begin{figure}[!htbp]
    \centering

    \begin{subfigure}[t]{0.96\textwidth}
        \centering
        \appAimage{Figures/Anexo_01/original_group_attributes.png}
        \caption{Initial table obtained from the active result file.}
        \label{fig:appendix_original_group_attributes}
    \end{subfigure}

    \vspace{0.8em}

    \begin{subfigure}[t]{0.96\textwidth}
        \centering
        \appAimage{Figures/Anexo_01/final_group_attributes.png}
        \caption{Final time-dependent and multi-case table.}
        \label{fig:appendix_final_group_attributes}
    \end{subfigure}

    \caption{Evolution of the group-attribute inspection area.}
    \label{fig:appendix_group_attributes_evolution}
\end{figure}

\FloatBarrier

\section{Linear and Radiative Conductors}
\label{appsec:conductor_modules_evolution}

The initial conductor tabs inspect the GL and GR matrices of the active case
and provide tables between selected groups. Their table values are tied to the
first stored conductor row, even when the source file contains a transient
solution. The final modules expose an output-time selector and preserve the
source--target sign convention in the generated table. Each case reports the
number of detected couplings and the corresponding total GL or GR value, and
the result can be exported as a separate table.

% GL capture: source OBC_BOX, destination RADIATOR, t = 0 s, generated table.
% Keep the count and total visible in the final image.
\begin{figure}[!htbp]
    \centering

    \begin{subfigure}[t]{0.96\textwidth}
        \centering
        \appAimage{Figures/Anexo_01/original_gl_conductors.png}
        \caption{Initial GL-conductor table.}
        \label{fig:appendix_original_gl_conductors}
    \end{subfigure}

    \vspace{0.8em}

    \begin{subfigure}[t]{0.96\textwidth}
        \centering
        \appAimage{Figures/Anexo_01/final_gl_conductors.png}
        \caption{Final GL-conductor table at the selected output time.}
        \label{fig:appendix_final_gl_conductors}
    \end{subfigure}

    \caption{Evolution of the linear-conductor inspection area.}
    \label{fig:appendix_gl_conductors_evolution}
\end{figure}

\FloatBarrier

% GR capture: source OBC_BOX, destination RADIATOR, t = 0 s, generated table.
% Use the same vertical crop as the GL comparison.
\begin{figure}[!htbp]
    \centering

    \begin{subfigure}[t]{0.96\textwidth}
        \centering
        \appAimage{Figures/Anexo_01/original_gr_conductors.png}
        \caption{Initial GR-conductor table.}
        \label{fig:appendix_original_gr_conductors}
    \end{subfigure}

    \vspace{0.8em}

    \begin{subfigure}[t]{0.96\textwidth}
        \centering
        \appAimage{Figures/Anexo_01/final_gr_conductors.png}
        \caption{Final GR-conductor table at the selected output time.}
        \label{fig:appendix_final_gr_conductors}
    \end{subfigure}

    \caption{Evolution of the radiative-conductor inspection area.}
    \label{fig:appendix_gr_conductors_evolution}
\end{figure}

\FloatBarrier

\section{Nodes and Group Study}
\label{appsec:nodes_group_study_evolution}

The nodal table in the initial application already supports attribute and
group selection, but filtering is applied directly while text is entered. The
final implementation adds explicit \textit{Apply filter} and
\textit{Clear filter} controls, records the selected output time and generates
case-specific or difference tables without modifying the imported data.

% Select BATTERY and OBC_BOX, attributes T, QI, QR and L, t = 0 s. Apply a
% label filter in the final capture so that the new Apply/Clear controls and
% active-filter message are visible.
\begin{figure}[!htbp]
    \centering

    \begin{subfigure}[t]{0.96\textwidth}
        \centering
        \appAimage{Figures/Anexo_01/original_nodes.png}
        \caption{Initial nodal table and direct text filter.}
        \label{fig:appendix_original_nodes}
    \end{subfigure}

    \vspace{0.8em}

    \begin{subfigure}[t]{0.96\textwidth}
        \centering
        \appAimage{Figures/Anexo_01/final_nodes.png}
        \caption{Final nodal table with explicit filter control.}
        \label{fig:appendix_final_nodes}
    \end{subfigure}

    \caption{Evolution of the nodal-inspection interface.}
    \label{fig:appendix_nodes_evolution}
\end{figure}

\FloatBarrier

The initial \textit{Group study} tab combines a group summary with internal
and external conductor inspection. The final tab retains this scope while
applying the common case selector, output-time control, representation mode
and export logic. GL and GR results are stored independently so that the
configuration associated with each conductor type can be restored without
altering the other analysis.

% Select BATTERY, t = 0 s, generate the summary and keep the internal/external
% GL and GR analysis controls visible. Use the same selected group in both
% versions.
\begin{figure}[!htbp]
    \centering

    \begin{subfigure}[t]{0.96\textwidth}
        \centering
        \appAimage{Figures/Anexo_01/original_group_study.png}
        \caption{Initial component summary and conductor inspection.}
        \label{fig:appendix_original_group_study}
    \end{subfigure}

    \vspace{0.8em}

    \begin{subfigure}[t]{0.96\textwidth}
        \centering
        \appAimage{Figures/Anexo_01/final_group_study.png}
        \caption{Final multi-case group and conductor analysis.}
        \label{fig:appendix_final_group_study}
    \end{subfigure}

    \caption{Evolution of the group-study interface.}
    \label{fig:appendix_group_study_evolution}
\end{figure}

\FloatBarrier

\section{Summary of the Inherited Modules}
\label{appsec:inherited_modules_summary}

Table~\ref{tab:appendix_inherited_module_summary} condenses the adaptations
shown in the preceding figures. The initial function of each module is
retained; the changes mainly affect case organisation, transient evaluation,
selection consistency and reproducibility.

% File: Tables/table_appendix_inherited_module_summary.tex
\begingroup
\small
\renewcommand{\arraystretch}{1.66}
\setlength{\tabcolsep}{3.5pt}
\begin{longtable}{|>{\raggedright\arraybackslash}p{0.15\textwidth}|>{\raggedright\arraybackslash}p{0.27\textwidth}|>{\raggedright\arraybackslash}p{0.47\textwidth}|}
    \caption{Adaptation of the inherited inspection modules.}
    \label{tab:appendix_inherited_module_summary}\\
    \hline
    \textbf{Module} & \textbf{Initial function} & \textbf{Final adaptation} \\
    \hline
    \endfirsthead

    \multicolumn{3}{c}{\tablename\ \thetable\ -- continued}\\
    \hline
    \textbf{Module} & \textbf{Initial function} & \textbf{Final adaptation} \\
    \hline
    \endhead

    \hline
    \endfoot

    Flow Explorer &
    Hierarchical heat-flow inspection for one active dataset. &
    Hierarchy-ordered selection, removal of redundant descendants, several
    cases, selected output time, overlay or difference representation, export
    and restoration. \\
    \hline
    Case information &
    Structural and coupling data from the active model. &
    Independent summaries, visible case names and level-based coupling
    inspection at a selected time. \\
    \hline
    Group attributes &
    Attribute table from the first or only stored result row. &
    Time-dependent tables for several cases and
    reference-minus-comparison results. \\
    \hline
    GL conductors &
    Conductive-coupling inspection for the active case. &
    Signed source--target table, selected time, coupling count, total GL,
    export and restoration. \\
    \hline
    GR conductors &
    Radiative-coupling inspection for the active case. &
    Signed source--target table, selected time, coupling count, total GR,
    export and restoration. \\
    \hline
    Nodes &
    Nodal table with immediate text filtering. &
    Multi-case and time-dependent tables, difference mode, explicit filtering
    and export. \\
    \hline
    Group study &
    Component summary and internal or external conductor inspection. &
    Common case and time controls, case-specific tables and independently
    restorable GL and GR analyses. \\
    \hline
\end{longtable}
\endgroup


\FloatBarrier
